\chapter{Experimento preliminar}
\label{ch:experimento}

El experimento preliminar valida el flujo de medición antes de la caracterización completa. Usa rutinas dominadas por una categoría conocida para comprobar que la plataforma muestra diferencias de potencia.

Reutiliza la cadena JTAG y la cadena de medición de potencia ya descritas, sin introducir una arquitectura de prueba adicional. 

% -----------------------------------------------------------------------
\section{Objetivo del experimento}
\label{sec:objetivo_experimento}

Verificar que rutinas dominadas por distintas categorías producen diferencias medibles y repetibles en la potencia promedio. 

% -----------------------------------------------------------------------
\section{Configuración de la prueba}
\label{sec:configuracion_experimento}

El SoC PULPissimo se sintetizará sobre la FPGA Nexys~A7-100T y ejecutará firmware bare-metal cargado mediante GDB, OpenOCD y FT232H/JTAG. La corriente se medirá con el INA240A1 y el ADS1115--LilyGO ESP32.

Cada rutina se ejecutará durante 10 a 60 segundos y el firmware marcará la ventana de medición con un pulso GPIO. La potencia promedio se calculará con la ecuación~\ref{eq:potencia_promedio}.

% -----------------------------------------------------------------------
\section{Rutinas de bucles dominados por categoría}
\label{sec:bucles_homogeneos_preliminares}

Las rutinas de prueba ejecutan un bloque reducido de instrucciones de una sola categoría. El objetivo es que esa categoría domine la ejecución, aunque existan instrucciones auxiliares de control.

% -----------------------------------------------------------------------
\section{Datos a recolectar}
\label{sec:datos_experimento}

Para cada rutina se registrarán la categoría dominante, el tiempo de adquisición, la potencia promedio y la dispersión entre repeticiones.

% -----------------------------------------------------------------------
\section{Criterio de viabilidad}
\label{sec:criterio_viabilidad}

El experimento será útil si rutinas de categorías distintas producen diferencias de potencia mayores que la variabilidad entre repeticiones. Si eso ocurre, queda justificada la caracterización posterior.

%%% Local Variables:
%%% mode: latex
%%% TeX-master: "main"
%%% End:
